\input{preamble}

\section*{Sommerfeld model}

Sommerfeld in 1916 derived the following energy eigenvalues for hydrogen.
These eigenvalues explain fine structure lines in the hydrogen spectrum.
\begin{equation*}
E=mc^2\left[1+\left(
\frac{\alpha}{n-k+\sqrt{k^2-\alpha^2}}\right)^2\right]^{-1/2}
\end{equation*}

Parameters $n$ and $k$ are the quantum numbers
\begin{align*}
n&=1,2,3,\ldots
\\
k&=1,2,\ldots,n
\end{align*}

Expand $E$ as a Taylor series in $\alpha$.
\begin{equation*}
E=mc^2\left[1-\frac{\alpha^2}{2n^2}-\frac{\alpha^4}{2kn^3}+\frac{3\alpha^4}{8n^4}
+\mathcal O(\alpha^6)\right]
\tag{1}
\end{equation*}

The leading term is electron rest-mass energy which can be discarded as unobservable in atomic spectra.
$\mathcal O(\alpha^6)$ is also discarded to obtain the approximation
\begin{equation*}
E=mc^2\left[-\frac{\alpha^2}{2n^2}-\frac{\alpha^4}{2kn^3}+\frac{3\alpha^4}{8n^4}\right]
\end{equation*}

Typically this is written as
\begin{equation*}
E=-\frac{\alpha^2mc^2}{2n^2}
\left[1+\frac{\alpha^2}{n}\left(\frac{1}{k}-\frac{3}{4n}\right)\right]
\tag{2}
\end{equation*}

By the substitution
\begin{equation*}
\alpha^2=\frac{e^4}{(4\pi\varepsilon_0\hbar c)^2}
\end{equation*}

we also have
\begin{equation*}
E=-\frac{e^4}{(4\pi\varepsilon_0\hbar)^2}
\frac{m}{2n^2}
\left[1+\frac{\alpha^2Z^2}{n}\left(\frac{1}{k}-\frac{3}{4n}\right)\right]
\end{equation*}

\href{https://georgeweigt.github.io/examples/sommerfeld-model-demo.html}{Eigenmath script}

\end{document}
